\documentclass[11pt]{article}
\title{Testing Stage-Dependent Injury-to-Epidermis Coupling in Xenopus laevis Tail Regeneration: An Unvalidated Mechanistic Hypothesis}
\date{}
\author{SIGMA AI-assisted research synthesis}
\usepackage[T1]{fontenc}
\usepackage[utf8]{inputenc}
\usepackage[margin=1in]{geometry}
\usepackage{amsmath,amssymb,booktabs,longtable}
\usepackage{microtype}
\usepackage{titlesec}
\emergencystretch=2em
\raggedbottom
\linespread{1.07}
\usepackage{xcolor}
\definecolor{blue}{HTML}{24485E}
\titleformat{\section}{\Large\bfseries\sffamily\color{blue}}{\thesection}{0.8em}{}
\titleformat{\subsection}{\large\bfseries\sffamily}{\thesubsection}{0.7em}{}
\setlength{\parskip}{.6em}
\setlength{\parindent}{0pt}
\begin{document}
\maketitle

\begin{abstract}
The developmental loss and subsequent recovery of tail-regeneration competence in Xenopus laevis motivates a distinction between failure to initiate a regenerative epidermal state and failure to couple injury responses to that state. This paper would develop a conditional hypothesis: early wound-response intensity is more strongly associated with subsequent regeneration-initiating-cell and extracellular-matrix program activation at stage $40$ than at stage $46$, after accounting for tissue composition and developmental differences. The proposed analysis would compare this interaction against stage-only, RIC/ECM-only, and established signaling-module baselines, using biological replicates rather than cells as independent observations.

However, the supplied evidence contains research proposals and explicit reports of absent completed results, not verified sample manifests, expression analyses, effect estimates, or independent biological validation. The appropriate contribution is therefore a bounded hypothesis and falsification plan, not a newly discovered predictor or explanation of the developmental switch. Novelty, quantitative performance, and causal relevance remain unresolved.
\end{abstract}
\section*{Evidence coverage limitation}
The evidence index omitted 1863 material saved research field fragment(s) across 6 source(s). These fragments were not reviewed for this paper. Interpretations that could depend on omitted material remain provisional; the complete frozen source snapshot and omission ledger are retained with the generation record.
\section{Introduction and biological context}
The biological question concerns variation in regeneration competence among X. laevis tadpoles, not merely separation of developmental-stage labels. The project describes a transient refractory interval and later recovery, but the proposed stage-$40$ versus stage-$46$ comparison does not itself establish how competence returns. An article should therefore distinguish a hypothesis about entry into refractoriness from an explanation of the complete loss-and-recovery trajectory. [C2, C11]

The supplied background identifies regeneration-initiating cells, their basal-epidermal origin, and ECM-remodeling defects as prior findings rather than candidate discoveries. The potentially discriminating question is whether early injury responses retain their association with later epidermal activation across developmental stages, rather than whether familiar regenerative programs are expressed. These background statements require confirmation from the underlying publications before they can serve as a cited literature synthesis. [C3, C4, C10]

Four explanations must remain distinguishable: failure of the epidermal transition, failure of injury-to-regeneration coupling, metabolic or tissue-composition gating, and analytical or developmental confounding. Nutritional observations attributed to X. tropicalis should motivate an alternative explanation, not establish the corresponding mechanism in X. laevis. The paper should preserve these alternatives until a conditional association and an independent intervention distinguish them. [C2, C7, C9]

\section{Methods and evidence}
Proposed evidence acquisition would begin with a sample-level manifest for GSE245320 and the nominated component series: bulk RNA-seq GSE245317/GSE245318, single-cell RNA-seq GSE245312, and spatial data GSE245313. The manifest must verify species, developmental stage, sampling time, tissue, animal or pool identity, biological replication, library preparation, batch, accessible files, and any regeneration outcomes. The reported bulk time points—$0$, $0.5$, and $6$ hours and $3$ days post-amputation—are starting metadata claims, not verified analytical coverage. [C5]

Proposed inference would use the independently sampled animal or documented biological pool as the unit of replication; cells, spatial locations, and technical libraries would remain nested measurements. Before testing temporal coupling, the analysis must establish whether early and later measurements share a defensible biological grouping. If they do not, the study may describe stage-specific population trajectories but must not manufacture paired observations or estimate animal-level temporal coupling from unrelated samples. [C6]

Subject to identifiability, the proposed comparison would add a stage-by-early-response interaction to simple baselines and quantify its incremental explanatory or predictive value. Module definitions, normalization, time windows, negative controls, multiplicity handling, and replicate-level resampling would be specified before evaluation; cell composition, batch, and stage-associated expression would be examined as competing explanations. No model has been fitted in the supplied evidence, and no numerical performance or uncertainty estimate is available. [C4, C7, C8]

Independent testing would combine a new phenotype-linked biological cohort with reciprocal rescue-and-block experiments, once an actionable coupling component is identified. The proposed design would include stage matching, sham and injury controls, appropriate treatment controls, toxicity assessment, target-engagement measurements, and blinded regeneration scoring at a prespecified endpoint. Neither a specific validated perturbation nor an executed experiment is supplied. [C9, C10]

\section{Biological scope and competing explanations}
The central mechanistic distinction is whether reduced tail-regeneration competence in Xenopus laevis reflects failure to initiate a regenerative epidermal state or failure to connect an earlier injury response to that state. These explanations could produce similar downstream deficiencies but locate the proposed developmental defect at different steps. The question is therefore not simply whether regeneration-associated programs differ between developmental stages, but what evidence would distinguish impaired epidermal activation from impaired injury-to-epidermis coupling.

Developmental-stage differences must also be separated from variation among animals. The motivating biological premise includes both a temporary reduction in regeneration competence and heterogeneous regeneration outcomes; a stage label should therefore not be treated as an individual animal’s outcome. The nominated stage $40$ versus stage $46$ comparison addresses a putatively competent–refractory contrast, not the complete loss-and-recovery sequence. No verified recovered-stage cohort is supplied, so this contrast alone cannot establish a mechanism for the later return of competence.

The supplied background attributes transient regeneration-initiating cells (RICs) derived from basal epidermal cells, together with extracellular-matrix (ECM) remodeling defects associated with regeneration failure, to a 2024 Genome Biology study. These observations define the intended prior-art boundary: identifying RICs or an ECM-associated deficiency again would not, by itself, establish the proposed contribution. Their detailed primary-source support remains to be confirmed from the publication, and the supplied material contains no completed expression analysis, predictive result, or mechanistic validation.

The appropriate contribution is thus a hypothesis-and-test plan; the absence of completed results is not evidence that any biological hypothesis has failed.

Under the epidermal-transition explanation, the relevant defect would lie in induction of the transient RIC/ECM-remodeling state. Under the coupling explanation, the distinguishing defect would instead concern the relationship between an early injury response and subsequent activation of that state. Reduced downstream activation alone would not discriminate these accounts: both permit that observation. Evidence favoring coupling failure would need to show that the dependence of later epidermal activation on the earlier response differs by developmental stage, rather than merely documenting a lower downstream signal in refractory animals.

Metabolic or tissue-composition gating provides a separate alternative to a defect in signal transfer. Under these accounts, an apparent association between early injury responses and later regenerative programs could depend on physiological context or on the abundance of relevant cell populations, rather than a changed response within the epidermal cells themselves. Distinguishing these possibilities would require evidence that separates changes in cell state from changes in tissue composition and assesses metabolic context rather than assuming it.

Nutritional explanations motivated by X. tropicalis would remain hypotheses for X. laevis: the supplied background cautions that cross-species stage behavior is context-dependent and that published accounts are not fully consistent.

Analytical and developmental confounding must remain alternatives to all of these biological explanations. A molecular contrast aligned with stage or batch would not, without further discrimination, identify an injury-specific defect; likewise, a rule that separates stage classes would not establish prediction of regeneration among animals within a stage. The mechanistic burden is therefore more specific than detecting a competent–refractory difference: it is to distinguish a change in the injury-to-epidermis relationship from downstream transition failure, physiological gating, and developmental differences.

This distinction shifts the argument from the presence or absence of known regenerative programs to the conditional relationship that would have to change for coupling failure to explain loss of competence.

\section{A conditional injury-to-epidermis coupling hypothesis}
The proposed injury-to-epidermis coupling hypothesis locates the developmental difference in the relationship between an early wound response and subsequent regenerative epidermal activation, rather than in either response alone. In Xenopus laevis tadpoles, the specific prediction is that stronger early injury responses are more positively associated with later regeneration-initiating-cell (RIC) and extracellular-matrix (ECM) program activation at stage $40$ than at stage $46$, conditional on biological and technical covariates. The supplied background identifies basal-epidermal RIC formation and ECM-remodeling defects as previously reported observations.

Their role here is to define a candidate downstream transition, not to constitute the proposed novelty.

The candidate early window is $0.5$ hours post-amputation, with the unamputated or immediate post-amputation reference at $0$ hours used, where the verified design permits, to distinguish injury-associated induction from pre-existing developmental differences. The candidate subsequent window is $6$ hours post-amputation. Early wound-response intensity and later RIC/ECM activation would require biologically justified module definitions rather than gene sets selected because they maximize the stage contrast. Measurements at $3$ days should be examined separately: the hypothesis does not assume that they index the same transition as the candidate early epidermal response.

These temporal assignments remain contingent on verified sample coverage and evidence that the proposed modules represent the intended processes at those times.

The directional commitment can be expressed by defining $\beta_{40}$ and $\beta_{46}$ as the conditional associations between early wound-response intensity and subsequent epidermal activation in the respective stages, for an estimand supported by the sampling design. The hypothesis predicts both $\beta_{40}>0$ and $\beta_{40}>\beta_{46}$. It does not require a negative association or complete absence of coupling at stage $46$; attenuated positive coupling would also satisfy the proposed direction. No numerical magnitude, threshold, or uncertainty interval is supplied by the available evidence, so these inequalities specify a testable prediction rather than an estimated result.

The interaction is essential because a difference in mean RIC/ECM activation between stages would not, by itself, establish altered coupling. Lower downstream activation could coexist with an unchanged relationship to early injury intensity, or could reflect a difference in the early response itself. Conversely, similar average injury responses could be accompanied by different subsequent activation relationships. The proposed signature therefore concerns whether developmental stage modifies the association between the two responses, beyond information available from stage-only, RIC/ECM-only, and established signaling-module baselines. A stage contrast alone would not establish prediction of regeneration outcomes among individual animals.

This formulation distinguishes an insufficient early response from a failure to translate an early response into the candidate epidermal program. If stage $46$ exhibited a weaker early response but retained the same conditional relationship to later activation as stage $40$, that pattern would argue against the specified coupling defect. If appreciable early responses were instead followed by disproportionately weak epidermal activation at stage $46$, the pattern would be consistent with impaired coupling. It would not uniquely identify that mechanism, however: an intrinsic inability to enter the downstream epidermal state could generate a similar association.

The hypothesis thus sharpens the comparison between mechanisms without treating temporal association as evidence of causal transfer.

The proposed downstream locus is the wound epidermis, not undifferentiated whole-tail transcription. RIC-related state activation and ECM-remodeling activity should be evaluated as related candidate readouts rather than assumed to be interchangeable. A bulk association that primarily tracked the abundance of RIC-like cells, another tissue component, or generalized healing would support a composition or healing explanation instead of the specified state-transition interpretation.

Spatial and single-cell observations could help distinguish these possibilities, but the supplied description indicates that the single-cell data predominantly characterize regenerating animals rather than matched snapshots across both stages; these modalities also belong to the same study and do not provide independent biological replication.

The available material therefore supports a precise hypothesis but supplies neither a measured coupling difference nor evidence that the relationship is previously unreported. Its substantive contribution is to change the mechanistic target from the presence of a known regenerative epidermal program to the stage-dependent relationship between injury intensity and that program. Evidence for that relationship would need to add information beyond stage, familiar signaling responses, and tissue composition; otherwise, the coupling formulation would not improve on the existing alternatives.

\section{Identifiability, baselines, and alternative explanations}
The proposed reanalysis must distinguish a stage-dependent temporal association from separate developmental differences in injury response and epidermal activation. That distinction depends first on whether the sampling design identifies a relationship between early and later measurements, rather than merely supplying measurements at both times. The supplied evidence establishes neither a verified sample manifest nor a numerical interaction effect, uncertainty interval, baseline advantage, or animal-level regeneration prediction; the analysis below therefore specifies the estimands and comparisons that a verified design could support.

The sample manifest should resolve the relationship between GSE245320 and its nominated bulk RNA-seq components, GSE245317 and GSE245318, before expression contrasts are defined. For each sample, it should establish species, developmental stage, post-amputation time, sampled tissue, biological replicate, pooling status, library identity, batch, expression units, and any linked regeneration phenotype. The reported stage $40$ versus stage $46$ comparison and collection times of $0$, $0.5$, and $6$ hours post-amputation and $3$ days are starting information, not a verified, fully replicated experimental design.

This distinction is consequential: absent combinations, uncertain pooling, or confounding between stage and batch would change which contrasts are estimable.

Biological replication must determine the unit of inference. Cells, spatial locations, or technical libraries derived from the same biological sample should not be counted as independent animals, and a pooled sample should not be treated as multiple independently measured individuals. Likewise, the bulk, single-cell, and spatial components belong to a shared study; agreement between them could strengthen interpretation within that study but would not constitute independent external replication. Whether the biological replicate is an individual tadpole or a documented pool must be established rather than inferred from the number of available libraries.

The decisive identifiability question is whether an early wound-response measurement has a defensible link to a later RIC/ECM measurement. Early and late samples from different animals cannot be paired by replicate number, similarity of expression, or common developmental stage and then analyzed as repeated observations of an animal. Such pairing would impose the very temporal relationship being tested. A documented relationship at another level, such as a replicated cohort, would instead define a cohort-level estimand and would not justify an animal-level interpretation.

If the samples are unmatched, the defensible analysis is a comparison of population trajectories: stage-specific changes from the corresponding pre-amputation reference and stage-by-time differences in module activation. These contrasts could establish whether early injury responses and later RIC/ECM programs differ between stages, but not whether animals with stronger early responses subsequently activate the epidermal program more strongly. In particular, ordering stage-level means in time does not identify a within-stage coupling slope. If no replicated linkage supports that slope, the temporal-coupling analysis should stop at this boundary rather than substitute an association between group averages.

If a suitable linkage is verified, a simple conditional model would make the coupling estimand explicit: \$R\_i=\textbackslash{}alpha+\textbackslash{}beta\_E E\_i+\textbackslash{}beta\_S S\_i+\textbackslash{}beta\_\{ES\}E\_iS\_i+\textbackslash{}boldsymbol\{\textbackslash{}gamma\}\textasciicircum{}\{\textbackslash{}mathsf\{T\}\}\textbackslash{}mathbf\{Z\}\_i+\textbackslash{}varepsilon\_i\$. Here, $R_i$ denotes later RIC/ECM activation, $E_i$ the linked early wound-response score, $S_i$ an indicator for stage $46$ relative to stage $40$, and $\mathbf{Z}_i$ a limited set of identifiable batch and composition covariates. The index $i$ must refer to the verified biological unit, not an arbitrarily paired sample; the proposed weaker association at stage $46$ corresponds to $\beta_{ES}<0$.

Interpreting this interaction would also require adequate within-stage variation and overlap in early-response scores, because near-separation by stage would make a slope difference difficult to distinguish from a developmental contrast.

Baseline comparisons must use the same endpoint and biological units. For later RIC/ECM activation, the interaction model should be compared with a stage-only model, an additive model without the interaction, and models based on established injury-signaling modules; a later RIC/ECM score cannot simultaneously serve as its own predictor. A RIC/ECM-only predictive baseline becomes meaningful for a distinct endpoint, such as a documented regeneration phenotype, only if that endpoint is genuinely available and linked to the relevant samples.

Under that condition, stage-only, RIC/ECM-only, and established-module baselines would test whether the proposed coupling signature contributes information beyond the known correlates rather than merely repackaging them.

Module definitions, candidate time contrasts, and the primary interaction should be specified before assessing predictive performance. Any outcome-dependent gene selection or tuning should be confined to training biological units, with all measurements from the same unit kept together during evaluation; held-out cells from a training animal would not supply an independent test. Uncertainty should be calculated at the biological-replicate level, with multiplicity control for additional modules or time windows and explicit inspection of collinearity and sensitivity to individual replicates.

If replication is insufficient for stable estimation or held-out evaluation, a fitted coefficient alone would not establish a reproducible quantitative rule.

Tissue composition is a competing explanation to examine, not merely a nuisance term to remove automatically. The proposed use of scRNA-seq GSE245312 and spatial GSE245313 is to ask whether a bulk signature reflects altered abundance of RIC-like cells, altered transcription within an epidermal state, or a broader wound-associated tissue mixture. Because the nominated single-cell data predominantly characterize regenerating animals rather than matched single-cell samples from both conditions, they should not be treated as direct evidence of a stage-specific cell-state difference.

Where coverage permits, replicate-aware cell-type summaries, marker-based composition estimates, and spatial localization could provide complementary checks on the interpretation of the bulk signal.

These checks should discriminate explanations without making composition adjustment synonymous with mechanistic proof. A signature that is accounted for by cellular abundance, or that lacks the proposed wound-epidermal localization, would weaken its interpretation as an injury-to-epidermis coupling measure. Conversely, a stage-dependent association that remains within comparable tissue or cell-state contexts would be more informative than a bulk association alone, but would still not establish causal transfer between the programs. A lower mean RIC/ECM response without a supported slope difference would remain compatible with failure of the downstream transition rather than specifically favoring coupling failure.

The evidentiary consequence is a sharper separation between developmental trajectory, conditional molecular association, and regeneration prediction. The nominated datasets can support the first only to the extent that their stage-by-time replication is verified, the second only if early and later observations have an identifiable biological relationship, and the third only with linked phenotypes and evaluation on held-out biological units. Thus, a difference between stages cannot by itself advance the central argument from impaired epidermal activation to impaired injury-to-epidermis coupling; that advance requires information not supplied by the stage contrast alone.

\section{Independent falsification and the limits of the proposed mechanism}
Reciprocal perturbation would test whether injury-to-epidermis coupling is a causal bottleneck in Xenopus laevis tail regeneration, rather than merely an association accompanying developmental stage. The discriminating question is whether selectively restoring the transfer from an early injury response to a later epidermal program can overcome stage-$46$ failure, and whether interrupting that same transfer can impair stage-$40$ regeneration. These tests would distinguish a coupling defect from a downstream epidermal-transition defect only if the intervention leaves the downstream machinery available for an independent assessment.

Under the coupling-defect hypothesis, restoration at stage $46$ should increase subsequent epidermal regeneration-initiating-cell (RIC) and extracellular-matrix (ECM) program activation and improve tail regrowth relative to matched injured controls. Conversely, selective interruption at stage $40$ should reduce both downstream activation and regrowth despite preservation of the initiating injury response. The reciprocal design therefore requires more than opposite changes in regeneration: it requires evidence that both interventions act on the nominated transfer step, with the predicted temporal ordering between target engagement, epidermal activation, and phenotype.

A downstream transition that remains inaccessible despite adequate injury input makes a different rescue prediction: restoring transfer alone need not restore the epidermal program or regrowth. However, failure to rescue would count against the proposed coupling bottleneck only if the relevant target were demonstrably engaged at the appropriate time and tissue. Likewise, directly forcing a downstream RIC/ECM program could bypass either an upstream coupling defect or a defect in normal transition induction; successful bypass would not, by itself, identify the endogenous lesion. Reciprocal blocking would strengthen the argument but would still require excluding nonspecific impairment of injury responses or tissue viability.

The supplied evidence identifies no selective coupling target, reagent, dose, or validated intervention window, so this remains an experimental logic rather than an executable rescue-and-block protocol. The candidate early response around $0.5$ hours post-amputation and subsequent activation around $6$ hours provide prospective temporal anchors, not established treatment boundaries. Before choosing an intervention, the proposed transfer step would need a measurable readout distinguishable from both injury intensity and the downstream program used to judge rescue. Otherwise, an increase in epidermal activation could be mistaken for proof of target engagement, making the mechanistic test circular.

Controls would need to separate a coupling-specific effect from altered injury severity, toxicity, and developmental mismatch. Each stage would require matched amputated controls receiving the corresponding control treatment, together with appropriate sham controls to assess effects not contingent on amputation. Toxicity and preservation of the early injury response should be assessed alongside localized epidermal activation rather than inferred from eventual regrowth alone.

A quantitative regeneration endpoint, such as a blinded morphological score at a prospectively specified post-amputation time, would need to be defined before testing, with biological replicates as the units of inference and sample-size assumptions tied to a stated effect of interest rather than an unsupported numerical expectation.

The joint molecular and morphological outcomes would determine how strongly the experiment bears on the mechanism. Restoration of epidermal activation without improved regrowth would argue that repairing the proposed transfer is insufficient to explain regeneration failure, even if it corrects an early molecular defect. Improved regrowth without the predicted epidermal response would challenge the proposed mediating role of that response or the adequacy of its measurement. A reciprocal loss of both readouts at stage $40$ would be informative only when toxicity and disruption of the initiating wound response have been excluded; otherwise, it would not distinguish selective uncoupling from generalized injury impairment.

Independent validation should use previously unseen X. laevis biological specimens with regeneration phenotypes, without selecting the validation cohort because it already exhibits the predicted relationship. The nominated bulk, single-cell, and spatial datasets share a study and cannot substitute for independent biological replication. To support prediction of variable outcomes among animals, molecular measurements would need a defensible link to the corresponding phenotypes; separating developmental stages alone would not meet that requirement.

No numerical interaction effect, uncertainty interval, advantage over the proposed baselines, or animal-level predictive performance is established in the supplied evidence, and the proposed prior-art assessment does not yet establish novelty.

The later recovery of competence is a separate boundary on the causal argument. A stage-$40$ versus stage-$46$ rescue-and-block comparison could test a candidate bottleneck in the nominated contrast, but no verified recovered-stage cohort or molecular trajectory is supplied. Extending the hypothesis to recovery would require testing whether the injury-to-epidermis dependency returns in phenotype-linked X. laevis specimens; context-dependent observations in X. tropicalis would not establish that trajectory. The evidence thus makes selective restoration of transfer, rather than generic enhancement of regeneration, the decisive causal question and limits any eventual explanation to the developmental transitions actually tested.

\section{Limitations}
\begin{itemize}
\item The selected evidence documents no completed research result. This is an evidence-availability limitation, not evidence against temporal coupling or any competing biological mechanism. [C1]
\item Accession contents, biological replication, longitudinal or cohort linkage, batch structure, and individual outcome labels have not been verified. These determine whether the proposed estimand is identifiable. [C5, C6]
\item The proposed interaction has no numerical effect estimate, uncertainty interval, or demonstrated advantage over known-mechanism baselines. [C8]
\item Shared-study modalities cannot provide independent external replication, and incomplete stage coverage in single-cell data may prevent the desired cell-state comparison. [C5, C7]
\item No completed prior-art ledger or independent perturbation evidence supports a novelty or mechanistic claim. The later recovery of competence also lies outside the verified scope of the proposed two-stage contrast. [C10, C11]
\end{itemize}
\section{Open obligations}
\begin{itemize}
\item Retrieve the underlying publications and build a dated, cited prior-art ledger that explicitly checks whether the proposed conditional interaction has already been reported. [C3, C10]
\item Verify the nominated GEO sample manifests, accessible files, biological units, pooling, batches, time points, and any phenotype linkage before defining the analysis population. [C5]
\item Determine whether the sampling design supports temporal coupling at an animal or defensible biological-group level; otherwise replace that estimand with explicitly population-level stage-by-time contrasts. [C6]
\item Prespecify module membership, timing, tissue localization criteria, baseline models, negative controls, multiplicity procedures, evaluation splits, and replicate-aware uncertainty estimation. [C4, C7]
\item Estimate the interaction and its incremental value numerically, or report that it is nonidentifiable, unstable, or unsupported rather than inventing a result. [C8]
\item Identify a selective perturbable coupling component, validate target engagement, and finalize intervention timing, treatment controls, blinded outcome scoring, sample-size assumptions, and ethical requirements. [C9]
\item Test the locked prediction in previously unseen phenotype-linked specimens or independently generated perturbation data. [C10]
\item Identify and test an independently verified recovered-competence condition before extending the hypothesis to the complete loss-and-recovery switch. [C11]
\end{itemize}
\section{Conclusions}
The supplied evidence does not establish why X. laevis tadpoles temporarily lose tail-regeneration competence. It supports a clearly bounded question about stage-dependent injury-to-epidermis coupling, not a discovered biological explanation. [C1, C2, C4]

The first decisive requirement is identifiability: the available biological units must support the proposed temporal association, and any interaction must add information beyond stage, known regenerative programs, and tissue composition. [C5, C6, C7]

A new mechanism requires traceable novelty, quantitative uncertainty, and independent biological confirmation. Until those obligations are met, the appropriate outcome is an unvalidated hypothesis and a proposed falsification strategy; later recovery remains an additional test. [C8, C9, C10, C11]

\section*{Authorship note}
AI-assisted synthesis from saved research; not independently peer reviewed.
\end{document}
